
\section{Method}

\subsection{Overview}

HALO-ReGLA runs two attention streams in parallel and combines them with 
per-head gates:

\begin{enumerate}
\item \textbf{Dense stream} $O^\text{dense}$: standard softmax attention.
\item \textbf{Linear stream} $O^\text{lin}$: rectified kernel attention.
\item \textbf{Gate} $g$: per-head per-token sigmoid.
\item \textbf{Fusion} $O = g \odot O^\text{dense} + (1 - g) \odot O^\text{lin}$.

\end{enumerate}
\textbf{Figure 4.1} illustrates the pipeline.

\begin{lstlisting}
        Q, K, V $\in$ R^{B $\times$ H $\times$ L $\times$ d_h}
                 |
        +--------+--------+
                           
   +---------+      +----------+
   | Dense   |      | Linear   |
   | softmax |      | (ELU+1)  |
   |   O_d   |      |   O_l    |
   +----+----+      +----+-----+
        |                |
        +--------+-------+
                  
        +-----------------+
        |  Gate g(x)      |
        |  g O_d+(1-g) O_l|
        +--------+--------+
                  
              out_proj
\end{lstlisting}

\subsection{Dense Stream}

Standard softmax attention with causal masking and padding:

\[
S_{ij} = \frac{Q_i^\top K_j}{\sqrt{d_h}}, \quad
S_{ij} \leftarrow S_{ij} + \text{mask}_{ij}, \quad
O^\text{dense}_i = \sum_j \text{softmax}(S_i)_j V_j
\]

\textbf{Complexity.} $\mathcal{O}(L^2 d_h)$ time, $\mathcal{O}(L^2)$ memory.

\textbf{Masks applied.}
\begin{itemize}
\item Causal: $S_{ij} = -\infty$ for $j > i$.
\item Padding: $S_{ij} = -\infty$ for invalid keys $j$.

\end{itemize}
\subsection{Linear Stream}

Rectified kernel attention with cumulative sums for causality:

\subsubsection{Feature Map}

\[
\phi(x) = \mathrm{elu}(x) + 1 \in \mathbb{R}^{d_h}
\]

Since $\mathrm{elu}(x) > -1$ for all $x$, we have $\phi(x) > 0$.

\subsubsection{Causal Formulation}

For each position $i$:

\[
\text{cum}^{(i)} = \sum_{j \le i} \phi(K_j) V_j^\top \in \mathbb{R}^{d_h \times d_h}
\]

\[
\text{den}^{(i)} = \sum_{j \le i} \phi(K_j) \in \mathbb{R}^{d_h}
\]

\[
O^\text{lin}_i = \frac{\phi(Q_i) \cdot \text{cum}^{(i)}}{\phi(Q_i) \cdot \text{den}^{(i)} + \epsilon}
\]

where $\epsilon = 10^{-5}$ prevents division by zero.

\subsubsection{Non-Causal Formulation}

For encoders (bidirectional), we use the same formulas without cumsum:

\[
O^\text{lin}_i = \frac{\phi(Q_i) \cdot \sum_j \phi(K_j) V_j^\top}{\phi(Q_i) \cdot \sum_j \phi(K_j) + \epsilon}
\]

\subsubsection{Padding Handling}

We zero out padding keys before the cumsum:

\[
\phi(K_j) \leftarrow \phi(K_j) \cdot \text{valid}_j, \quad
V_j \leftarrow V_j \cdot \text{valid}_j
\]

This ensures padding contributes nothing to numerator or denominator.

\subsubsection{Numerical Stability}

The cumsum is computed in FP32 to avoid overflow:

\[
\text{cum}^{(i)} = \text{cumsum}(\phi(K) V^\top).\text{float}()
\]

We cast back to the model dtype after cumsum.

\subsection{Fusion Gate}

\subsubsection{Parameterization}

For each head $h$, a linear projection computes the gate logit:

\[
\text{logit}_{h,i} = w_h^\top x_i + b_h
\]

where $x_i \in \mathbb{R}^d$ is the input hidden state at position $i$. The 
gate is:

\[
g_{h,i} = \sigma(\text{logit}_{h,i}) \in (0, 1)
\]

\subsubsection{Output}

\[
O_{h,i} = g_{h,i} \cdot O^\text{dense}_{h,i} + (1 - g_{h,i}) \cdot O^\text{lin}_{h,i}
\]

\subsubsection{Initialization}

We initialize $b_h = 1.8$, so $\sigma(1.8) \approx 0.858$ at initialization. 
This biases the model toward the \textbf{dense stream}, which matches the 
pretrained behavior.

\textbf{Table 4.1: Gate values at various biases.}

\begin{table}[H]
\centering
\small
\begin{tabular}{lll}
\toprule
\textbf{$b_h$} & \textbf{$\sigma(b_h)$} & \textbf{Interpretation} \\
\midrule
$-2$ & 0.12 \& Mostly linear \\
$0$ & 0.50 \& Balanced \\
$1.8$ & 0.86 & \textbf{Mostly dense (default)} \\
$3$ & 0.95 \& Almost purely dense \\
\bottomrule
\end{tabular}
\end{table}

\subsubsection{Per-Head vs Global Gate}

We use \textbf{per-head} gates. This allows different heads to prefer different 
streams, which matches empirical observations that some heads are local 
(better served by linear) and others are global (better served by dense).

\subsection{Complete Forward Pass}

\textbf{Algorithm 1:} HALO-ReGLA forward
\begin{lstlisting}
Input:  hidden_states $\in$ R^{B $\times$ L $\times$ d}
        attention_mask, position_embeddings
        is_causal
Output: output $\in$ R^{B $\times$ L $\times$ d}

 1: Q, K, V $\leftarrow$ proj(hidden_states)
 2: Q, K $\leftarrow$ apply_rope(Q, K, position_embeddings)
 3: K, V $\leftarrow$ repeat_kv(K, V, H / H_kv) if GQA
 4:
 5: # Dense stream
 6: S $\leftarrow$ Q K^T / sqrt(d_h)
 7: S $\leftarrow$ S + causal_mask if is_causal
 8: S $\leftarrow$ S + pad_mask if attention_mask
 9: A_dense $\leftarrow$ softmax(S)
10: O_dense $\leftarrow$ A_dense @ V
11:
12: # Linear stream
13: Q_l $\leftarrow$ elu(Q) + 1
14: K_l $\leftarrow$ elu(K) + 1
15: K_l, V_mask $\leftarrow$ apply_pad_mask(K_l, V, attention_mask)
16:
17: if is_causal:
18:     cum_kv $\leftarrow$ cumsum(K_l   V_mask)         # (B, H, L, d_h, d_h)
19:     num $\leftarrow$ (Q_l.unsqueeze(-2) @ cum_kv).squeeze(-2)
20:     cum_k $\leftarrow$ cumsum(K_l)
21:     den $\leftarrow$ (Q_l * cum_k).sum(-1, keepdim=True)
22: else:
23:     KV $\leftarrow$ K_l^T @ V_mask
24:     Z $\leftarrow$ K_l.sum(dim=2, keepdim=True)
25:     num $\leftarrow$ Q_l @ KV
26:     den $\leftarrow$ Q_l @ Z^T
27: O_lin $\leftarrow$ num / den.clamp(min=1e-5)
28: O_lin $\leftarrow$ O_lin * valid_mask
29:
30: # Gate and fuse
31: g $\leftarrow$ sigmoid(gate(hidden_states))          # (B, H, L, 1)
32: O $\leftarrow$ g * O_dense + (1 - g) * O_lin
33:
34: O $\leftarrow$ O.reshape(B, L, H * d_h)
35: O $\leftarrow$ out_proj(O)
36: return O
\end{lstlisting}

\subsection{Complexity Analysis}

\begin{table}[H]
\centering
\small
\begin{tabular}{llll}
\toprule
\textbf{Component} & \textbf{Dense} & \textbf{Linear} & \textbf{HALO-ReGLA} \\
\midrule
Time & $\mathcal{O}(L^2 d_h)$ & $\mathcal{O}(L d_h^2)$ & $\mathcal{O}(L^2 d_h + L d_h^2)$ \\
Memory & $\mathcal{O}(L^2)$ & $\mathcal{O}(d_h^2)$ & $\mathcal{O}(L^2)$ \\
Causal cumsum & -- & $\mathcal{O}(L d_h^2)$ \& Same \\
\bottomrule
\end{tabular}
\end{table}

\textbf{Note.} HALO-ReGLA does \textbf{not} reduce total complexity because the dense 
stream is still computed. The benefit is in the \textbf{gate}: if the gate 
converges to 0 for many heads, we can drop those heads' dense computations 
post-training.

To realize speedup, one would:
\begin{enumerate}
\item Train with both streams.
\item Identify heads where $\bar{g}_h < \theta$ (e.g., 0.1).
\item Convert those heads to pure linear inference.

\end{enumerate}
\subsection{Parameter Analysis}

\begin{table}[H]
\centering
\small
\begin{tabular}{lll}
\toprule
\textbf{Component} & \textbf{Baseline} & \textbf{HALO-ReGLA} \\
\midrule
$W_Q$ & $d H d_h$ \& Same \\
$W_K$ & $d H_{kv} d_h$ \& Same \\
$W_V$ & $d H_{kv} d_h$ \& Same \\
$W_O$ & $H d_h d$ \& Same \\
\textbf{Gate} & 0 & $d H + H$ \\
LoRA & 0 & $2 r (d H d_h + \cdots)$ \\
\bottomrule
\end{tabular}
\end{table}

The gate adds $d H + H$ parameters (one linear per head). For $d = 768$, 
$H = 12$: \textbf{9,216 parameters}. Negligible.

\subsection{Theoretical Notes}

\subsubsection{Approximation of Dense by Linear}

\textbf{Claim.} If the attention distribution is smooth (no sharp peaks), the 
linear approximation is close to dense.

\textbf{Sketch.} Linear attention approximates softmax with kernel $\phi$. For 
bounded inputs and $d_h \to \infty$, PORF gives unbiased approximation. Our 
ELU kernel is heuristic but works well in practice.

\subsubsection{Gate Convergence}

\textbf{Claim.} The gate converges to a stationary point where:

\[
\frac{\partial \mathcal{L}}{\partial g_h} = 0 \Rightarrow O^\text{dense}_h = O^\text{lin}_h
\]

or the gate is at a boundary ($g_h \to 0$ or $g_h \to 1$).

\textbf{Proof.} At optimum, $\partial \mathcal{L} / \partial g = (\nabla_O \mathcal{L}) \cdot (O^\text{dense} - O^\text{lin}) = 0$. 

\subsubsection{Causality Guarantee}

The causal prefix-sum formulation ensures position $i$ only attends to 
$j \le i$. This is because $\text{cum}^{(i)}$ depends only on $\{\phi(K\_j), 
V\_j : j \le i\}$. $\square$
\begin{lstlisting}

---
\end{lstlisting}